\documentclass[12pt,oneside]{book}

% ==================== METADATOS ====================
\newcommand{\universidad}{Universidad de Buenos Aires (UBA)}
\newcommand{\facultad}{Facultad de Derecho}
\newcommand{\materia}{Contratos Civiles y Comerciales}
\newcommand{\titulo}{El contrato de compraventa inmobiliaria: título y modo suficiente, instrumentos previos y boleto}
\newcommand{\autor}{Isaac Edgar Camacho Ocampo}
\newcommand{\anio}{2026}

% ==================== PAQUETES ====================
\usepackage[utf8]{inputenc}
\usepackage[T1]{fontenc}
\usepackage[spanish,es-nodecimaldot]{babel}
\usepackage[
    top=2.5cm,
    bottom=2.5cm,
    left=3cm,
    right=2.5cm,
    headheight=15pt
]{geometry}
\usepackage{amsmath}
\usepackage{amssymb}
\usepackage{mathtools}
\usepackage{graphicx}
\usepackage{float}
\usepackage{caption}
\usepackage{subcaption}
\usepackage{booktabs}
\usepackage{array}
\usepackage{longtable}
\usepackage{tabularx}
\usepackage{enumitem}
\usepackage{xcolor}
\usepackage[most]{tcolorbox}
\usepackage{fancyhdr}
\usepackage{ragged2e}
\usepackage[
    colorlinks=true,
    linkcolor=blue!50!black,
    citecolor=green!40!black,
    urlcolor=blue!60!black,
    pdftitle={\titulo},
    pdfauthor={\autor},
    pdfsubject={Apuntes universitarios},
    bookmarks=true
]{hyperref}

% ==================== CONFIGURACIÓN ====================
\graphicspath{
    {./img/}
    {./graficos/}
    {./figuras/}
}
\setcounter{tocdepth}{2}
\setlength{\parindent}{0pt}
\setlength{\parskip}{0.5em}
\setlist[itemize]{
    topsep=0.3em,
    itemsep=0.2em
}
\setlist[enumerate]{
    topsep=0.3em,
    itemsep=0.2em
}
\clubpenalty=10000
\widowpenalty=10000

% ==================== COMANDOS ====================
\newcommand{\abs}[1]{\left|#1\right|}
\newcommand{\norm}[1]{\left\lVert#1\right\rVert}
\newcommand{\vect}[1]{\mathbf{#1}}
\newcommand{\deriv}[2]{\frac{d#1}{d#2}}
\newcommand{\pderiv}[2]{\frac{\partial#1}{\partial#2}}

% ==================== ENTORNOS / CAJAS ====================
\newtcolorbox{definition}[1]{
    enhanced,
    breakable,
    title={Definición: #1},
    fonttitle=\bfseries,
    colback=blue!3,
    colframe=blue!50!black,
    boxrule=0.8pt,
    arc=2mm
}

\newtcolorbox{important}{
    enhanced,
    breakable,
    title={Importante},
    fonttitle=\bfseries,
    colback=yellow!8,
    colframe=orange!70!black,
    boxrule=0.8pt,
    arc=2mm
}

\newtcolorbox{example}[1]{
    enhanced,
    breakable,
    title={Ejemplo: #1},
    fonttitle=\bfseries,
    colback=green!4,
    colframe=green!40!black,
    boxrule=0.8pt,
    arc=2mm
}

\newtcolorbox{formula}[1]{
    enhanced,
    breakable,
    title={Fórmula / Regla: #1},
    fonttitle=\bfseries,
    colback=purple!4,
    colframe=purple!50!black,
    boxrule=0.8pt,
    arc=2mm
}

\newtcolorbox{exercise}[1]{
    enhanced,
    breakable,
    title={Ejercicio: #1},
    fonttitle=\bfseries,
    colback=gray!5,
    colframe=gray!60!black,
    boxrule=0.8pt,
    arc=2mm
}

\newtcolorbox{note}{
    enhanced,
    breakable,
    title={Nota},
    fonttitle=\bfseries,
    colback=white,
    colframe=gray!50,
    boxrule=0.6pt,
    arc=2mm
}

% ==================== CONFIGURACIÓN FINAL ====================
\pagestyle{fancy}
\fancyhf{}
\fancyhead[L]{\nouppercase{\leftmark}}
\fancyhead[R]{\materia}
\fancyfoot[C]{\thepage}
\renewcommand{\headrulewidth}{0.4pt}
\renewcommand{\footrulewidth}{0pt}
\fancypagestyle{plain}{
    \fancyhf{}
    \fancyfoot[C]{\thepage}
    \renewcommand{\headrulewidth}{0pt}
}

\renewcommand{\arraystretch}{1.25}

% ==================== DOCUMENTO ====================
\begin{document}

% ---------- PORTADA ----------
\begin{titlepage}

\thispagestyle{empty}

\begin{center}

\vspace*{1cm}

{\large \textsc{\universidad}}

\vspace{0.2cm}

{\large \textsc{\facultad}}

\vspace{3cm}

{\Huge \textbf{\materia}}

\vspace{1cm}

{\LARGE \titulo}

\vspace{0.8cm}

\textit{Comprender los conceptos es el primer paso para convertir el estudio en conocimiento.}

\vspace{1.2cm}

\rule{0.70\textwidth}{0.5pt}

\vspace{0.5cm}

{\large Apuntes de estudio}

\vspace{0.5cm}

\rule{0.70\textwidth}{0.5pt}

\vfill

{\large \autor}

\vspace{0.3cm}

{\large \anio}

\end{center}

\vfill

\begin{flushleft}
\textit{Este es un modesto aporte a los estudiantes de ``\materia'' no pretende sustituir el material oficial de la materia.}
\end{flushleft}

\end{titlepage}

% ---------- ÍNDICE ----------
\tableofcontents

% ---------- PANORAMA ----------
\chapter*{Panorama de la clase}
\addcontentsline{toc}{chapter}{Panorama de la clase}
\markboth{Panorama de la clase}{Panorama de la clase}

\textbf{Compraventa: título, modo y boleto.} Clase del 26 de octubre de 2026.

La clase parte de una pregunta práctica: ¿cuándo una persona pasa a ser dueña de un inmueble que compró? Para adquirir el dominio por un acto entre vivos hacen falta dos piezas de igual valor: el \textbf{título suficiente} (con requisitos de fondo y de forma) y el \textbf{modo suficiente} (la tradición posesoria). Luego se recorre, como en una operación real, la cadena de instrumentos que se firman antes de la escritura: corretaje, promesas de venta y de compra, reserva, seña y refuerzo, con especial atención a la diferencia entre reserva y seña. Por último se presenta el boleto de compraventa, que el Código regula sin definir, y las dos posturas doctrinarias sobre si es un contrato preliminar o el contrato definitivo.

\section*{Cómo se relacionan los temas}

Los tres bloques de la clase son tres preguntas sobre una misma compra (por ejemplo, la de un departamento):

\begin{itemize}
    \item \textbf{Bloque 1:} la meta. ¿Cuándo soy realmente dueño? (título más modo).
    \item \textbf{Bloque 2:} el camino. ¿Qué papeles se firman mientras tanto? (corretaje, promesas, reserva, seña, refuerzo).
    \item \textbf{Bloque 3:} un punto intermedio muy discutido. Cuando se firma el boleto, ¿ya se compró o todavía no?
\end{itemize}

\begin{table}[H]
\centering
\small
\begin{tabularx}{\textwidth}{@{}>{\raggedright\arraybackslash}p{0.18\textwidth}>{\raggedright\arraybackslash}X@{}}
\toprule
\textbf{Etapa} & \textbf{Qué ocurre} \\
\midrule
1. Se publica la casa & Corretaje: el dueño autoriza a una inmobiliaria a ofrecerla. \\
2. Se hace una oferta & Promesa de compra: se propone comprar bajo ciertas condiciones. \\
3. Reserva & Poca plata para que retiren el inmueble del mercado mientras se negocia. \\
4. Seña y refuerzo & Con la oferta aceptada, se entrega más dinero para dar firmeza. \\
5. Boleto & Mientras se hace el estudio de títulos: ¿ya se compró o todavía no? \\
6. Escritura y llaves & Escritura (título) más entrega de la posesión (modo): recién ahí se es dueño. \\
\bottomrule
\end{tabularx}
\end{table}

\section*{Utilidad de los contenidos}

\begin{itemize}
    \item \textbf{Redactar y asesorar sin errores.} Si se pide una «seña» y lo que corresponde es una reserva (o al revés), cambia por completo el dinero que se pierde o se devuelve doblado si alguien se arrepiente. Importa el contenido del instrumento y no su título.
    \item \textbf{Definir la estrategia de un reclamo.} Si el boleto es un contrato preliminar, vence al año y solo permite reclamar daños; si es el contrato definitivo, se puede exigir la escrituración.
    \item \textbf{Proteger al cliente frente a terceros.} Los artículos 1170 y 1171 fijan cuándo el boleto prevalece sobre un embargo o resiste una quiebra: buena fe, fecha cierta, publicidad suficiente y el 25\,\% del precio pagado.
    \item \textbf{Controlar la validez del título.} Capacidad, legitimación, inhabilidades del artículo 1002 y escritura pública son la lista de verificación de quien revise una operación.
    \item \textbf{En la vida personal.} Comprar o vender una vivienda suele ser la decisión económica más grande de una vida: permite entender qué se firma en cada etapa y qué riesgo se asume (por ejemplo, que la plata entregada como seña, salvo pacto en contrario, cuenta como principio de pago y confirma el acto; que nadie puede vender un derecho que no tiene; que ser dueño exige escritura más posesión real; y que la operación lleva como mínimo unos 30 días de trámites).
\end{itemize}

\begin{note}
Este material es de estudio y no reemplaza el asesoramiento de un profesional matriculado para un caso concreto.
\end{note}

% =====================================================================
\chapter[Título y modo suficiente]{Título y modo suficiente para adquirir el dominio}

\section{Adquisición derivada del dominio}

El dominio se adquiere en virtud de un contrato, pero el contrato no es la única causa posible. Se puede adquirir por distintas causas, incluso tratándose de derechos reales derivados, por actos entre vivos, que se ejercen por la posesión. Otros ejemplos son la \textbf{donación} y la \textbf{dación en pago}: no existe solo el contrato de compraventa.

\begin{note}
El fragmento inicial de la clase está incompleto: la grabación comienza con la clase ya iniciada.
\end{note}

Ubicados en el contrato de compraventa, se analizan los requisitos que debe reunir el \textbf{título} para ser suficiente y los que debe reunir el \textbf{modo} para ser suficiente.

\section{Título suficiente: requisitos de fondo y de forma}

El título, para ser suficiente, tiene requisitos de \textbf{fondo} y de \textbf{forma}.

\begin{itemize}
    \item \textbf{Fondo:} legitimación y capacidad.
    \item \textbf{Forma:} el modo de exteriorizar la voluntad.
\end{itemize}

\section{Legitimación}

\begin{definition}{Legitimación}
La legitimación tiene que ver con el principio \textbf{nemo plus iuris}: nadie puede transmitir un derecho mejor ni más extenso que el que tiene. Quien dispone debe ser titular del derecho que pretende disponer.
\end{definition}

\section{Capacidad e inhabilidades para contratar}

\subsection{Capacidad de ejercicio}

Se distingue la capacidad de ejercicio de la capacidad de derecho. La \textbf{capacidad de ejercicio} consiste en contar con capacidad para otorgar el acto, y \textbf{se presume}: hoy todos somos capaces y, solo excepcionalmente, la capacidad se restringe para determinados actos.

\begin{itemize}
    \item La restricción es judicial y excepcional, y se dispone para actos específicos. El juez indica las \textbf{medidas de apoyo} necesarias para asistir a la persona en el otorgamiento de esos actos: el apoyo \emph{asiste}, no \emph{reemplaza}.
    \item Que una persona tenga la capacidad restringida no significa que no pueda vender. Habrá que leer la sentencia, ver si la capacidad está restringida para ese acto y, si lo está, cumplir con la medida de apoyo indicada por el juez.
    \item En última instancia existe la \textbf{incapacidad total}: cuando la persona no puede interactuar con su entorno ni expresar su voluntad y el apoyo es insuficiente. Solo como último recurso se declara incapaz a la persona y se designa un curador. Tampoco esto significa que el incapaz no pueda vender: para vender bienes de un incapaz se requiere un requisito especial, la autorización judicial.
\end{itemize}

\subsection{Capacidad de derecho}

La \textbf{capacidad de derecho} consiste en ser titular de deberes y derechos, y tener la aptitud para disponer. Las incapacidades de derecho están enumeradas por la ley: son \textbf{taxativas}, enunciadas expresamente.

\begin{table}[H]
\centering
\small
\begin{tabularx}{\textwidth}{@{}>{\raggedright\arraybackslash}p{0.19\textwidth}>{\raggedright\arraybackslash}X>{\raggedright\arraybackslash}X@{}}
\toprule
 & \textbf{Capacidad de ejercicio} & \textbf{Capacidad de derecho} \\
\midrule
\textbf{Qué es} & Contar con capacidad para otorgar el acto. & Ser titular de deberes y derechos; tener aptitud para disponer. \\
\textbf{Regla general} & Se presume: hoy todos somos capaces. & Las incapacidades están enunciadas expresamente por la ley y son taxativas. \\
\textbf{Excepciones} & Restricción judicial excepcional, para actos específicos, con medidas de apoyo que asisten y no reemplazan. En última instancia, incapacidad total con curador. & Inhabilidades legales, como las del artículo 1002 del Código. \\
\textbf{Qué implica para vender} & Un capaz restringido o un incapaz no «no pueden vender»: hay que leer la sentencia, cumplir el apoyo o contar con autorización judicial. & Quien tiene la inhabilidad no puede comprar ni vender en esos supuestos. \\
\bottomrule
\end{tabularx}
\caption*{Comparación entre capacidad de ejercicio y capacidad de derecho.}
\end{table}

\subsection{Inhabilidades: el artículo 1002}

El artículo 1002 del Código enumera expresamente quiénes no pueden celebrar contratos: entre otros, los funcionarios públicos y los cónyuges entre sí. Los cónyuges casados bajo el régimen de comunidad no pueden contratar entre sí: marido y mujer no pueden venderse ni comprarse recíprocamente bienes, por incapacidad de derecho.

\begin{note}
\textbf{Artículo 1002 (Código Civil y Comercial). Inhabilidades especiales.} No pueden contratar en interés propio:
\begin{enumerate}[label=\alph*),topsep=0.2em]
    \item los funcionarios públicos, respecto de bienes de cuya administración o enajenación están o han estado encargados;
    \item los jueces, funcionarios y auxiliares de la justicia, los árbitros y mediadores, y sus auxiliares, respecto de bienes relacionados con procesos en los que intervienen o han intervenido;
    \item los abogados y procuradores, respecto de bienes litigiosos en procesos en los que intervienen o han intervenido;
    \item los cónyuges, bajo el régimen de comunidad, entre sí.
\end{enumerate}
Los albaceas que no son herederos no pueden celebrar contrato de compraventa sobre los bienes de las testamentarias que estén a su cargo.

\emph{Texto del Código incorporado como referencia en el apunte original; el docente lo enumera oralmente.}
\end{note}

En el ámbito judicial, los jueces y auxiliares de justicia no pueden contratar respecto de bienes de los procesos en que hayan intervenido: se trata de una cuestión de \textbf{imparcialidad}. Lo mismo rige para los funcionarios administrativos, respecto del proceso donde han intervenido, por la misma causa. La prohibición alcanza también a la actuación \textbf{por interpósita persona}: por ejemplo, el esposo de la jueza no puede comprar un bien de una sucesión que tramitó ante ella. Los abogados, por su parte, no pueden contratar respecto de bienes litigiosos del proceso donde hayan intervenido.

\section{La forma: escritura pública}

\textbf{Forma} es el modo de exteriorizar la voluntad; en el ámbito jurídico, el modo de expresarla. Se trata aquí de transmisiones derivadas, por actos entre vivos, de derechos reales que se ejercen por la posesión, lo que remite necesariamente a los \textbf{inmuebles}. Tratándose de inmuebles se exige \textbf{escritura pública}, conforme al artículo 1017. Solo se exceptúa un caso puntual: la \textbf{subasta}, que consiste en comprar un bien por un método de oferta pública al mejor postor. Escritura pública, entonces, con la excepción de la subasta.

\begin{note}
\textbf{Artículo 1017 (Código Civil y Comercial). Escritura pública.} Deben ser otorgados por escritura pública los contratos que tienen por objeto la adquisición, modificación o extinción de derechos reales sobre inmuebles. Quedan exceptuados los casos en que el acto es realizado mediante subasta proveniente de ejecución judicial o administrativa. \emph{[\ldots]}

\emph{Texto del Código incorporado como referencia (se transcribe el inciso a); el apunte original aclara que debe verificarse contra la edición oficial.}
\end{note}

\section{Justo título}

Cuando no se reúne alguno de los requisitos del título suficiente se plantea el \textbf{justo título}. El defecto que lo origina es un problema de \textbf{fondo}, no de forma. Por ejemplo: el poder estaba revocado. Un título así implica mala fe de alguna de las partes.
% TODO: contenido incompleto o ambiguo en las notas originales (pasaje poco claro en el audio)

\begin{example}{Error en el objeto}
Una persona se cree dueña de la unidad funcional cinco y en realidad ocupa desde siempre la cuatro. En ese caso no hay título: hay un error en el objeto, por más que crea que sí lo hay.
\end{example}

\section{Modo suficiente: la tradición posesoria}

Aun si se firma la escritura a favor del adquirente y esta reúne los requisitos de fondo y de forma, todavía no se adquiere el dominio: falta el \textbf{modo suficiente}. El modo tiene \textbf{igual valor} que el título; ambos se necesitan por igual y, si falta uno, no hay dominio. Uno sin el otro no sirve.

El modo suficiente es la \textbf{tradición posesoria traslativa de dominio}. La tradición, por sí sola, es solo un acto de entrega y recepción; la relación de poder que se busca es la \textbf{posesión}. No sirve cualquier tradición, sino la traslativa de dominio, porque hay otros derechos reales que también se ejercen con la posesión, como el usufructo.

\begin{definition}{Posesión y tenencia}
Entre las relaciones de poder se distinguen la posesión y la tenencia.
\begin{itemize}
    \item \textbf{Posesión:} se da cuando no se reconoce en otro un señorío superior al que el derecho ofrece.
    \item \textbf{Tenencia:} se da cuando, al contrario, se reconoce que otro tiene la posesión y se tiene solo la tenencia.
\end{itemize}
\end{definition}

\begin{table}[H]
\centering
\small
\begin{tabularx}{\textwidth}{@{}>{\raggedright\arraybackslash}p{0.22\textwidth}>{\raggedright\arraybackslash}X>{\raggedright\arraybackslash}X@{}}
\toprule
 & \textbf{Posesión} & \textbf{Tenencia} \\
\midrule
\textbf{Cuándo se da} & Cuando no se reconoce en otro un señorío superior al que el derecho ofrece. & Cuando se reconoce que otro tiene la posesión. \\
\textbf{Qué busca el comprador} & La tradición que le entregue la posesión, y no cualquiera: la traslativa de dominio. & No es lo que busca: con la tenencia no adquiere el dominio. \\
\bottomrule
\end{tabularx}
\end{table}

El modo suficiente consiste en \textbf{actos de entrega y recepción}, o actos de uno solo con el asentimiento del otro. Lo importante es que las \textbf{meras declaraciones} sobre la entrega de la posesión no tienen efectos frente a terceros: lo que importa es que efectivamente se reciba la posesión.

Esto significa que el adquirente pueda (es una facultad, no un deber) ejercer actos concretos y materiales sobre la cosa. No implica estar en contacto directo y permanente con la cosa, ni reformarla para que todos vean que se tiene la posesión. Basta con tener la posibilidad de ejercer actos concretos y materiales sobre la cosa, o que la cosa esté bajo la \textbf{custodia exclusiva} del adquirente, sin reconocer en otros un señorío superior.

Si en la escritura se omite, por ejemplo, consignar que el comprador se encuentra en posesión, que la recibió antes del acto o que no hay posesión de terceros (cláusulas que suelen incluirse), no ocurre nada, porque las meras declaraciones sobre darse por poseído o desposeído no tienen efectos frente a terceros. Lo importante frente a terceros es la posibilidad concreta de ejercer esos actos.

\begin{important}
La tradición no se prueba con una cláusula de la escritura. Lo decisivo frente a terceros es la posibilidad concreta de ejercer actos materiales sobre la cosa, o que esta quede bajo la custodia exclusiva del comprador.
\end{important}

\section{Recapitulación: título más modo}

\begin{formula}{Adquisición del dominio}
\textbf{Dominio} = título suficiente (fondo: legitimación y capacidad; forma: escritura pública) + modo suficiente (tradición posesoria traslativa de dominio). Ambos tienen igual valor; si falta uno, no hay dominio.
\end{formula}

Estos requisitos son para adquirir el derecho real de dominio.

% ---------- Cornell Cap. 1 ----------
\section{Cuadro Cornell: título y modo suficiente}

{\small
\begin{longtable}{@{}>{\raggedright\arraybackslash}p{0.28\textwidth}>{\raggedright\arraybackslash}p{0.68\textwidth}@{}}
\toprule
\textbf{Preguntas / palabras clave} & \textbf{Notas / respuestas} \\
\midrule
\endfirsthead
\toprule
\textbf{Preguntas / palabras clave} & \textbf{Notas / respuestas} \\
\midrule
\endhead
\bottomrule
\endfoot

\textbf{¿Cómo se adquiere el dominio por actos entre vivos?} &
Con título suficiente más modo suficiente, que tienen igual valor: si falta uno, no hay dominio. La compraventa es solo una de las causas posibles (también la donación o la dación en pago). \\[0.6em]

\textbf{¿Qué requisitos tiene el título suficiente?} &
De fondo: legitimación y capacidad. De forma: el modo de exteriorizar la voluntad. \\[0.6em]

\textbf{¿Qué es la legitimación?} &
Aplicación del principio \emph{nemo plus iuris}: nadie puede transmitir un derecho mejor ni más extenso que el que tiene. Hay que ser titular del derecho que se pretende disponer. \\[0.6em]

\textbf{¿Qué diferencia hay entre capacidad de ejercicio y de derecho?} &
La capacidad de ejercicio se presume; la restricción judicial es excepcional, para actos específicos y con medidas de apoyo que asisten y no reemplazan. La incapacidad total (con curador) es el último recurso y exige autorización judicial para vender. La capacidad de derecho es la aptitud para disponer; sus limitaciones son taxativas. \\[0.6em]

\textbf{¿Quiénes no pueden contratar según el art. 1002?} &
Funcionarios públicos, jueces y auxiliares de la justicia, abogados y procuradores, respecto de bienes de procesos en que intervinieron, y cónyuges en comunidad entre sí. Tampoco por interpósita persona (ejemplo: el esposo de la jueza). \\[0.6em]

\textbf{¿Qué forma exige la compraventa de inmuebles?} &
Escritura pública (art.~1017), con la excepción de la subasta (oferta pública al mejor postor). \\[0.6em]

\textbf{¿Qué es el justo título?} &
El título con un problema de fondo, como un poder revocado. Si hay error en el objeto (creer ser dueño de la unidad cinco y ocupar la cuatro), no hay título. \\[0.6em]

\textbf{¿Qué es el modo suficiente?} &
La tradición posesoria traslativa de dominio: actos de entrega y recepción. Posesión: no se reconoce en otro un señorío superior; tenencia: se reconoce que otro tiene la posesión. Importa la posibilidad concreta de ejercer actos materiales o la custodia exclusiva; las meras declaraciones de la escritura no valen frente a terceros. \\

\midrule
\multicolumn{2}{@{}p{\textwidth}@{}}{%
\textbf{Resumen:} Para adquirir el dominio por acto entre vivos hacen falta, con igual valor, un título suficiente (fondo: legitimación y capacidad; forma: escritura pública, salvo subasta) y un modo suficiente (tradición posesoria traslativa de dominio). Los defectos de fondo dan lugar al justo título; frente a terceros importa la posibilidad concreta de ejercer actos materiales sobre la cosa y no las declaraciones de la escritura.} \\
\end{longtable}
}

% =====================================================================
\chapter[Instrumentos previos a la escritura]{Instrumentos previos a la escritura}

\section{La operación inmobiliaria: el recorrido}

Para responder si el contrato de compraventa es formal o no formal se parte de lo ya visto: la definición del contrato de compraventa, sus efectos obligacionales y no traslativos de por sí, y lo necesario para adquirir el dominio. Para ello se enuncian todos los instrumentos que suelen firmarse con motivo de una operación inmobiliaria. El objetivo es identificar el contrato de compraventa, clasificarlo como formal o no formal e identificar la transmisión de dominio.

Entre los instrumentos \textbf{preliminares} a la firma de un contrato definitivo (como el boleto de compraventa) se mencionan: la \textbf{reserva}, la \textbf{seña} (planteada en clase como pregunta), el \textbf{contrato de corretaje} (autorización de venta) y las \textbf{promesas de venta} o \textbf{de compra}. Cuál se utiliza depende de quién hace la oferta.

\begin{important}
Todos estos instrumentos \textbf{no son obligatorios}. Suelen firmarse con motivo de una operación inmobiliaria que no se hace instantáneamente: es un proceso que lleva \textbf{como mínimo 30 días}, y en ese intervalo se suelen firmar distintos instrumentos, ninguno de los cuales es obligatorio.
\end{important}

\section{Contrato de corretaje}

\begin{definition}{Contrato de corretaje}
Es la autorización que dan los propietarios a la inmobiliaria o, mejor dicho, al profesional que se encarga de comercializar la propiedad. También se lo llama autorización de venta o autorización de comercialización. Consiste en una autorización que el titular de un bien da a un tercero para que lo ofrezca en el comercio a fin de celebrar un determinado acto.
\end{definition}

Hoy está regulado por el Código. En algún momento se discutió si era un contrato de corretaje o un mandato; hoy se entiende que es un contrato de corretaje.

Si la autorización es exclusivamente para que se venda (no para que se permute) y el corredor cumple el objetivo, es decir, encuentra un comprador dentro del plazo de vigencia (generalmente tiene uno), tiene derecho a su \textbf{retribución} aun cuando alguna de las partes luego se arrepienta y no se celebre el contrato. El corredor ya hizo su trabajo, que fue encontrar un comprador; si después el vendedor se arrepintió de vender, no es problema del corredor.

\section{Promesas de venta y de compra}

\begin{definition}{Promesa de venta}
Acto jurídico \textbf{unilateral} por el cual una persona dueña de un bien se obliga a vender, en ciertas condiciones, a un tercero. Si el tercero firma esta promesa, no significa que acepta el contrato, sino que lo va a estudiar y va a ver si acepta o no.
\end{definition}

\begin{definition}{Promesa de compra}
Lo mismo que la promesa de venta, pero a la inversa: una persona se obliga a comprar un bien en determinadas circunstancias y, si firma el titular, no se obliga a vender sino a estudiar la propuesta.
\end{definition}

En el contrato de corretaje el contrato es entre el corredor inmobiliario y el vendedor: \textbf{no hay contrato con el comprador}. En cambio, las promesas de venta y de compra son actos jurídicos unilaterales, mientras que el contrato de corretaje es bilateral. Hasta ese momento el comprador no participa.

\begin{table}[H]
\centering
\small
\begin{tabularx}{\textwidth}{@{}>{\raggedright\arraybackslash}p{0.17\textwidth}>{\raggedright\arraybackslash}X>{\raggedright\arraybackslash}X>{\raggedright\arraybackslash}X@{}}
\toprule
 & \textbf{Corretaje} & \textbf{Promesa de venta} & \textbf{Promesa de compra} \\
\midrule
\textbf{Naturaleza} & Contrato bilateral, regulado por el Código. & Acto jurídico unilateral. & Acto jurídico unilateral (inversa de la promesa de venta). \\
\textbf{Quiénes intervienen} & Corredor y vendedor. Todavía no participa el comprador. & El titular se obliga a vender a un tercero, en ciertas condiciones. & El comprador se obliga a comprar al titular, en ciertas condiciones. \\
\textbf{Qué compromete} & Autorización para ofrecer el bien; retribución del corredor si encuentra comprador en el plazo, aunque no se celebre el contrato. & El tercero que firma no acepta el contrato: lo va a estudiar. & El titular no se obliga a vender: estudia la propuesta. \\
\bottomrule
\end{tabularx}
\end{table}

\section{La reserva}

\begin{definition}{Reserva}
Contrato \textbf{innominado} (no regulado en el Código) y \textbf{preliminar}, cuyo objetivo específico es poder llegar a la celebración de un contrato definitivo, como el boleto de compraventa en el caso de la compraventa. Puede darse en cualquier operación inmobiliaria (permuta, locación), pero aquí se circunscribe a la compraventa.
\end{definition}

En este contrato el comprador entrega al vendedor una \textbf{pequeña suma de dinero} en concepto de que tiene interés en negociar las condiciones del contrato definitivo. A cambio, el vendedor se compromete a \textbf{retirar el inmueble del mercado}.

Todos estos instrumentos preliminares tienen como \textbf{plazo máximo de vigencia un año}: puede ser un año o menos. La reserva se estructura así: primero las partes, el objeto y, al hablar del precio, se dice que en ese acto se entrega cierta cantidad de dinero y que otra cantidad se entregará al suscribir el boleto de compraventa. Las partes pactan así su interés de llegar en algún momento a un acuerdo para celebrar un contrato definitivo.

La reserva es el acto jurídico unilateral en el que una persona interesada en comprar manifiesta su intención de obligarse a la compra de un bien determinado. Si el vendedor acepta, hay otro acto jurídico unilateral: la conformidad del vendedor. Si no lo acepta, puede presentar una contraoferta o contrapropuesta, con la cual también se tiene que seguir negociando. En todo caso, la reserva es un contrato preliminar.

\section{La seña (arras)}

\begin{definition}{Seña}
La seña \textbf{no es un contrato}, sino una \textbf{cláusula accidental} dentro de un contrato: no es obligatoria ni esencial, y las partes la agregan por propia voluntad. Está regulada en los artículos 1059 y 1060 del Código. También se la conoce como \textbf{arras} o \textbf{señal}.
\end{definition}

En la práctica pueden verse reservas tituladas como señas, pero son reservas. Si el título dice «Seña» cuando todavía se están negociando las condiciones del contrato definitivo, no hay seña: hay una reserva.

\begin{important}
Hay que verificar siempre el \textbf{contenido} del instrumento y no su título. Los efectos jurídicos de una reserva y de una seña son completamente diferentes, principalmente en lo económico (el dinero que se entrega) y en lo que ocurre en caso de incumplimiento.
\end{important}

\subsection{La seña: Código de Vélez y Código Civil y Comercial}

En el \textbf{Código de Vélez} la seña tenía \textbf{efecto penitencial}: aunque se entregaba dinero para asegurar el cumplimiento del contrato, las partes tenían la facultad de arrepentirse. Si el comprador se arrepentía, perdía automáticamente el dinero; si el vendedor se arrepentía, debía devolver la seña recibida más el doble, es decir, devolverla «doblada» (lo recibido más otro tanto).

Con la sanción del Código Civil y Comercial, en 2015, esto cambió.

\begin{note}
\textbf{Artículo 1059 (Código Civil y Comercial). Señal.} La entrega de señal o arras se interpreta como confirmatoria del acto, excepto que las partes convengan la facultad de arrepentirse. En tal caso, quien entregó la señal la pierde en beneficio de la otra parte, y quien la recibió debe restituirla doblada.

\emph{Texto leído en clase por una alumna.}
\end{note}

Hoy la seña tiene \textbf{efecto confirmatorio como regla}: quien la entrega se compromete a celebrar el contrato definitivo. Lo que debe pactarse expresamente es que las partes tengan la facultad de arrepentirse, es decir, que la seña sea \textbf{penitencial}.

El cambio es importante porque el dinero entregado sirve como \textbf{principio de ejecución} y como \textbf{cumplimiento parcial}.

\begin{example}{Seña confirmatoria}
Si se quiere comprar un inmueble, en un caso hipotético, de \$500 y se entregan \$100, restan \$400: el dinero entregado ya cuenta como principio de ejecución y como cumplimiento parcial.
\end{example}

Si las partes pactan que la seña sea penitencial, los efectos son los mismos que en el Código de Vélez: si se arrepiente el comprador, pierde la seña que dio; si se arrepiente el vendedor, debe devolver la seña recibida más otro tanto igual, es decir, doblada.

\begin{table}[H]
\centering
\small
\begin{tabularx}{\textwidth}{@{}>{\raggedright\arraybackslash}p{0.22\textwidth}>{\raggedright\arraybackslash}X>{\raggedright\arraybackslash}X@{}}
\toprule
 & \textbf{Código de Vélez} & \textbf{Código Civil y Comercial (art.~1059)} \\
\midrule
\textbf{Efecto de la seña} & Penitencial. & Confirmatorio como regla; penitencial solo si las partes lo pactan. \\
\textbf{Qué significa} & Las partes podían arrepentirse. & Confirma el acto: quien entrega la seña se compromete a celebrar el contrato definitivo. \\
\textbf{Si se arrepiente el comprador} & Pierde automáticamente el dinero entregado. & Si la seña es penitencial (pactada): pierde la seña que dio. \\
\textbf{Si se arrepiente el vendedor} & Devuelve lo recibido más otro tanto (doblado). & Si la seña es penitencial (pactada): devuelve lo recibido más otro tanto (doblado). \\
\textbf{Función del dinero} & Asegurar el cumplimiento del contrato. & Principio de ejecución y cumplimiento parcial. \\
\bottomrule
\end{tabularx}
\end{table}

Hasta aquí siempre se trata de \textbf{derechos personales}: todavía no hay constitución de ningún derecho real.

Las partes deben pactar lo relativo al incumplimiento. En las reservas y en las señas siempre se pacta qué ocurre en caso de incumplimiento de la parte compradora y de la parte vendedora, y eso debe dejarse expreso. Aunque no se diga que la seña tiene carácter penitencial, se entiende que, al pactarse el incumplimiento de cada parte, se le está dando ese efecto. Sin embargo, la regla hoy es que la seña es confirmatoria del acto: quien entrega el dinero ya tiene la intención de celebrar un contrato definitivo.

En la reserva, la suma que se entrega es muy baja y el acto es unilateral; si el vendedor no lo acepta, la inmobiliaria que interviene devuelve el dinero y no hay mayores efectos. En cambio, el dinero de la seña suele quedar \textbf{caucionado}, en la escribanía o en la inmobiliaria, hasta que salga bien el estudio de títulos o se reúnan todas las condiciones para escriturar (para lo cual se requieren varios certificados registrales y administrativos). Todo esto las partes lo deben prever.

\section{Reserva versus seña}

\begin{itemize}
    \item La \textbf{reserva} es el precontrato: el instrumento que contiene la voluntad del comprador de firmar un futuro contrato y de otorgar una suma de dinero como muestra de seriedad de la propuesta, a cambio de que el vendedor retire el inmueble del mercado.
    \item La \textbf{seña} no es un contrato per se, sino una cláusula accidental de un contrato en el que también se ofrece un dinero como parte de garantía del cumplimiento. La regla es que tiene carácter confirmatorio, a menos que las partes pacten que sea penitencial.
\end{itemize}

\begin{table}[H]
\centering
\small
\begin{tabularx}{\textwidth}{@{}>{\raggedright\arraybackslash}p{0.20\textwidth}>{\raggedright\arraybackslash}X>{\raggedright\arraybackslash}X@{}}
\toprule
 & \textbf{Reserva} & \textbf{Seña} \\
\midrule
\textbf{Naturaleza} & Contrato preliminar (precontrato), innominado; acto unilateral del interesado en comprar. & No es un contrato: es una cláusula accidental dentro de un contrato. \\
\textbf{Regulación} & No está regulada en el Código. & Artículos 1059 y 1060 del Código. También se la llama arras o señal. \\
\textbf{Función} & Expresar la voluntad de firmar un futuro contrato; el dinero muestra la seriedad de la propuesta. & El dinero es parte de la garantía del cumplimiento del contrato. \\
\textbf{Dinero} & Suma baja; si el vendedor no acepta, se devuelve. & Confirmatoria como regla; penitencial solo si se pacta. Suele quedar caucionada hasta poder escriturar. \\
\textbf{Qué hace el vendedor} & Retira el inmueble del mercado. & Según lo pactado para el caso de incumplimiento de cada parte. \\
\bottomrule
\end{tabularx}
\end{table}

\begin{note}
\textbf{Fallo recomendado:} \emph{Diluch, Pablo Fabio Armando c/ Más 3 S.R.L. y otros s/ daños y perjuicios.} Según indica el docente, «Más 3» es la inmobiliaria. El fallo define bien la seña y la reserva y permite ejemplificarlas con un caso concreto; se recomienda su lectura para la clase siguiente.
\end{note}

\section{El refuerzo}

El \textbf{refuerzo} se da cuando la oferta es aceptada y, para darle más firmeza a la operación, se entrega más dinero.

% ---------- Cornell Cap. 2 ----------
\section{Cuadro Cornell: instrumentos previos}

{\small
\begin{longtable}{@{}>{\raggedright\arraybackslash}p{0.28\textwidth}>{\raggedright\arraybackslash}p{0.68\textwidth}@{}}
\toprule
\textbf{Preguntas / palabras clave} & \textbf{Notas / respuestas} \\
\midrule
\endfirsthead
\toprule
\textbf{Preguntas / palabras clave} & \textbf{Notas / respuestas} \\
\midrule
\endhead
\bottomrule
\endfoot

\textbf{¿Qué se firma en una operación inmobiliaria?} &
Corretaje, promesas de venta y de compra, reserva, seña, refuerzo, boleto y escritura. Los instrumentos previos no son obligatorios, duran como máximo un año y el proceso lleva al menos 30 días. \\[0.6em]

\textbf{¿Cómo funcionan el corretaje y las promesas?} &
Corretaje: contrato bilateral entre corredor y vendedor, regulado; el corredor cobra si encuentra comprador en el plazo aunque no se celebre el contrato. Promesas: actos unilaterales; quien las recibe no se obliga, solo estudia la propuesta. \\[0.6em]

\textbf{¿Qué es la reserva?} &
Contrato preliminar innominado: el comprador entrega una suma baja para negociar el contrato definitivo y el vendedor retira el inmueble del mercado. Si el vendedor acepta hay conformidad; si no, puede contraofertar o se devuelve el dinero. \\[0.6em]

\textbf{¿Qué efecto tiene hoy la seña? (art.~1059)} &
Es una cláusula accidental. Por regla es confirmatoria (el dinero es principio de ejecución). Solo es penitencial si se pacta: quien la entregó la pierde, y quien la recibió la devuelve doblada. En el Código de Vélez era penitencial. \\[0.6em]

\textbf{¿Cómo se distinguen reserva y seña?} &
Verificando el contenido del instrumento y no su título. Tienen efectos económicos y de incumplimiento distintos. Fallo recomendado: Diluch c/ Más 3 S.R.L. \\[0.6em]

\textbf{¿Qué es el refuerzo?} &
Dinero adicional que se entrega cuando la oferta es aceptada, para darle más firmeza a la operación. \\

\midrule
\multicolumn{2}{@{}p{\textwidth}@{}}{%
\textbf{Resumen:} Antes de la escritura la operación recorre varios instrumentos no obligatorios (corretaje, promesas, reserva, seña y refuerzo) que no transmiten el dominio y generan solo derechos personales. La reserva (un precontrato con suma baja) y la seña (una cláusula accidental, hoy confirmatoria salvo pacto penitencial) se distinguen por su contenido y no por el título que lleven.} \\
\end{longtable}
}

% =====================================================================
\chapter[El boleto de compraventa]{El boleto de compraventa}

\section{El boleto en el Código: artículos 1170 y 1171}

El boleto de compraventa se puede firmar con motivo de una operación inmobiliaria y no es algo raro. Con él surge una inquietud: \textbf{no está definido por el Código}, pero dos artículos se refieren a sus consecuencias: el 1170 y el 1171.

\begin{note}
\textbf{Artículo 1170 (Código Civil y Comercial). Boleto de compraventa de inmuebles.} El derecho del comprador de buena fe tiene prioridad sobre el de terceros que hayan trabado cautelares sobre el inmueble vendido, si:
\begin{enumerate}[label=\alph*),topsep=0.2em]
    \item el comprador contrató con el titular registral o puede subrogarse en la posición jurídica de quien lo hizo mediante un perfecto eslabonamiento entre los sucesivos adquirentes;
    \item el comprador pagó como mínimo el 25\,\% del precio con anterioridad a la traba de la cautelar;
    \item el boleto tiene fecha cierta;
    \item la adquisición tiene publicidad suficiente, sea registral, sea posesoria.
\end{enumerate}
\emph{Texto leído en clase por un alumno. En la grabación el inciso d) se entiende «procesal»; el apunte original lo transcribe según el texto del Código.}
\end{note}

\begin{note}
\textbf{Artículo 1171 (Código Civil y Comercial). Oponibilidad del boleto en el concurso o quiebra.} Los boletos de compraventa de inmuebles de fecha cierta otorgados a favor de adquirentes de buena fe son oponibles al concurso o quiebra del vendedor si se hubiera abonado como mínimo el 25\,\% del precio. El juez debe disponer que se otorgue la respectiva escritura pública. El comprador puede cumplir sus obligaciones en el plazo convenido. En caso de que la prestación a cargo del comprador sea a plazo, debe constituirse hipoteca en primer grado sobre el bien en garantía del saldo de precio.

\emph{Texto leído en clase por un alumno.}
\end{note}

No hay duda de que el Código \textbf{no define} el boleto de compraventa: regula las consecuencias de haberlo firmado. Tampoco hay duda de que está inserto en la parte de compraventa (no en el capítulo de permuta), y de que se refiere a quien firma como adquirente o comprador; en algún inciso se refiere también a la adquisición del bien.

El Código nuevo podría haber dado una definición y terminado con la discusión sobre la naturaleza jurídica del boleto, pero no lo hizo: regula las consecuencias frente a terceros, pues habla de \textbf{oponibilidad}, es decir, de hacer valer algo frente a otros. Además, el boleto es un instrumento habitual y familiar, no extraordinario.

\begin{important}
El Código no define el boleto de compraventa: regula solo sus consecuencias frente a terceros (arts.~1170 y 1171). Lo ubica dentro del capítulo de compraventa y habla del «comprador», lo que alimenta el debate sobre su naturaleza.
\end{important}

\section{Naturaleza jurídica en debate: contrato preliminar o definitivo}

Como el Código no definió el boleto, se mantiene la discusión vigente durante el ordenamiento anterior: si se trata de un \textbf{contrato preliminar} o del \textbf{contrato definitivo} de compraventa.

\begin{definition}{Contrato preliminar}
Hoy el Código establece qué es un contrato preliminar: un acuerdo sobre los elementos esenciales de un futuro contrato. Como es un futuro contrato, no es el contrato mismo: las partes acuerdan los elementos esenciales de un contrato que celebrarán en algún momento.
\end{definition}

% TODO: contenido incompleto o ambiguo en las notas originales (pasaje con elipsis en la explicación oral)

\begin{note}
\textbf{Artículo 994 (Código Civil y Comercial). Disposiciones generales (contratos preliminares).} Los contratos preliminares deben contener el acuerdo sobre los elementos esenciales particulares que identifiquen el contrato futuro definitivo.

El plazo de vigencia de las promesas previstas en esta Sección es de un año, o el menor que convengan las partes, quienes pueden renovarlo a su vencimiento.

\emph{Texto del Código incorporado como referencia; el docente lo explica oralmente.}
\end{note}

Los contratos preliminares tienen un \textbf{plazo de vigencia máximo de un año}: es un plazo de caducidad, contado desde que se firma. Además, en caso de incumplimiento del acuerdo sobre el futuro contrato, solo se pueden reclamar \textbf{daños y perjuicios}.

Hay quienes consideran que el boleto es un contrato preliminar, porque entienden que no es el contrato definitivo sino el acuerdo sobre lo básico: cosa y precio. Hay otros que consideran que es el contrato definitivo. Hay argumentos de peso y referentes de renombre en ambos lados (entre ellos Borda en uno de los lados; el otro nombre no se distingue con claridad en el audio).
% TODO: contenido incompleto o ambiguo en las notas originales (nombre poco claro en el audio)

No se trata de una mera discusión doctrinaria para clasificar: las implicancias jurídicas de adoptar una u otra postura son muy diferentes. Si se adopta la postura de que es un contrato preliminar, sujeto a un plazo de caducidad de un año, a los dieciocho meses no se podría exigir la escrituración. Si el comprador no reclama dentro del año y no le cumplen, podría reclamar solo daños y perjuicios; el inmueble, pareciera que no.

\begin{important}
No es una discusión meramente teórica. Si el boleto es un contrato preliminar, vence al año y ante el incumplimiento solo se reclaman daños y perjuicios. Si es el contrato definitivo, se puede exigir la escrituración.
\end{important}

\section{Postura del contrato preliminar}

Esta postura entiende que el \textbf{boleto no es el contrato de compraventa}, porque los contratos de compraventa inmobiliaria requieren escritura pública, y en el boleto solo se definen la cosa que se va a vender y el precio que se va a pagar: los dos elementos esenciales. Para que el contrato de compraventa tenga lugar se requiere otra manifestación de voluntad. Entienden que con el boleto todavía no se celebra la compraventa, porque las partes no quieren celebrarla aún y por eso firman el boleto; por algún motivo no firman la escritura, que consideran la forma que debe tener ese contrato. Se basan en los artículos \textbf{994} y \textbf{1017}.

Leídos esos dos artículos (contrato preliminar como acuerdo sobre elementos esenciales; escritura pública para los contratos referidos a la transmisión o constitución de derechos reales sobre inmuebles) pareciera que esta postura tiene razón. Hay argumentos de peso de ambos lados. El problema es el plazo de caducidad y la posibilidad de reclamar solo daños.

\section{Postura contractualista y cierre de la clase}

La \textbf{postura contractualista} sostiene que el boleto de compraventa \textbf{es el contrato de compraventa en sí}. Según la definición que da el Código del contrato de compraventa, es el instrumento por el cual uno se obliga a transmitir el dominio y el otro se obliga a pagar un precio: la \textbf{asunción de la obligación} hace el contrato. Si del boleto surge la asunción de esas dos obligaciones, hay contrato de compraventa, y es un contrato definitivo.

Argumentos de esta postura:
\begin{itemize}
    \item \textbf{Libertad de forma.} El Código no exige una forma especial para el contrato de compraventa. Quienes sostienen la otra postura confunden el contrato de compraventa con la ejecución de la obligación principal que emerge del contrato: la transmisión de dominio, que sin duda exige escritura pública. Confunden el momento de celebración y nacimiento del contrato con el momento posterior de ejecución de la obligación que nace.
    \item \textbf{El contenido y no la forma.} Para analizar la naturaleza jurídica de un instrumento no hay que detenerse en la forma sino en lo que surge del instrumento mismo. El boleto de compraventa puede celebrarse incluso por escritura pública; es habitual con los poderes especiales irrevocables: se firma por escritura el boleto y se acompaña con un poder especial irrevocable. La naturaleza jurídica no puede mutar en función de la forma en que se exterioriza la voluntad: aunque el boleto se haga por escritura, todavía no hay transmisión de dominio ni se transmite ningún derecho real.
    \item \textbf{La parte especial.} El Código, al regular la compraventa, tampoco exige en la parte especial la escritura pública, como sí lo exige el capítulo específico de donación para el contrato de donación. Esto refuerza la argumentación.
\end{itemize}

\begin{table}[H]
\centering
\small
\begin{tabularx}{\textwidth}{@{}>{\raggedright\arraybackslash}p{0.19\textwidth}>{\raggedright\arraybackslash}X>{\raggedright\arraybackslash}X@{}}
\toprule
 & \textbf{Contrato preliminar} & \textbf{Contrato definitivo (contractualista)} \\
\midrule
\textbf{Qué es el boleto} & Un acuerdo sobre los elementos esenciales (cosa y precio) de un futuro contrato de compraventa. & El contrato de compraventa en sí. \\
\textbf{Argumento central} & La compraventa inmobiliaria requiere escritura pública; hace falta otra manifestación de voluntad para celebrarla. & Del boleto surgen las dos obligaciones que definen la compraventa: transmitir el dominio y pagar un precio. \\
\textbf{Normas en que se apoya} & Artículos 994 y 1017. & Definición del contrato de compraventa y principio general de libertad de forma; el capítulo de compraventa no exige escritura (a diferencia del de donación). \\
\textbf{Respuesta a la otra postura} & Se confunde el boleto con un acuerdo previo; para el contrato falta la forma. & La otra postura confunde la celebración del contrato con la ejecución de su obligación principal (la transmisión del dominio, que sí exige escritura). La naturaleza se define por el contenido y no por la forma. \\
\textbf{Consecuencia práctica} & Plazo de caducidad de un año; ante el incumplimiento, solo daños y perjuicios. & Permite exigir la escrituración: el comprador no queda limitado a reclamar daños. \\
\bottomrule
\end{tabularx}
\end{table}

La clase se interrumpe en este punto y el tema continúa en la clase siguiente, porque hay mucha información.

% ---------- Cornell Cap. 3 ----------
\section{Cuadro Cornell: el boleto de compraventa}

{\small
\begin{longtable}{@{}>{\raggedright\arraybackslash}p{0.28\textwidth}>{\raggedright\arraybackslash}p{0.68\textwidth}@{}}
\toprule
\textbf{Preguntas / palabras clave} & \textbf{Notas / respuestas} \\
\midrule
\endfirsthead
\toprule
\textbf{Preguntas / palabras clave} & \textbf{Notas / respuestas} \\
\midrule
\endhead
\bottomrule
\endfoot

\textbf{¿Qué dice el Código del boleto? (arts.~1170 y 1171)} &
No lo define; solo regula sus consecuencias frente a terceros. Prevalece sobre cautelares si hay buena fe, fecha cierta, publicidad suficiente, contrato con el titular registral y al menos el 25\,\% del precio pagado; es oponible al concurso o quiebra con fecha cierta, buena fe y 25\,\% abonado. \\[0.6em]

\textbf{¿Por qué importa si el boleto es preliminar o definitivo?} &
Si es preliminar (art.~994): vigencia máxima de un año, plazo de caducidad y solo se reclaman daños y perjuicios. Si es definitivo: se puede exigir la escrituración. \\[0.6em]

\textbf{¿Qué sostiene la postura del contrato preliminar?} &
El boleto solo acuerda cosa y precio; la compraventa inmobiliaria exige escritura pública, de modo que falta otra manifestación de voluntad. Se apoya en los artículos 994 y 1017. \\[0.6em]

\textbf{¿Qué sostiene la postura contractualista?} &
Del boleto surgen las dos obligaciones de la compraventa, y rige la libertad de forma. La otra postura confunde la celebración con la ejecución. La naturaleza se define por el contenido: el boleto puede ser por escritura y aun así no transmite dominio. \\

\midrule
\multicolumn{2}{@{}p{\textwidth}@{}}{%
\textbf{Resumen:} El boleto de compraventa, que el Código no define y solo regula en sus efectos frente a terceros (arts.~1170 y 1171), divide a la doctrina: para unos es un contrato preliminar (con plazo de caducidad de un año y reclamo limitado a daños), para otros es el contrato definitivo (con posibilidad de exigir la escrituración). Elegir una postura u otra cambia lo que se puede reclamar.} \\
\end{longtable}
}

% ---------- Síntesis final ----------
\chapter*{Síntesis final}
\addcontentsline{toc}{chapter}{Síntesis final}
\markboth{Síntesis final}{Síntesis final}

\begin{important}
Para ser dueño de un inmueble no alcanza con firmar un papel: hacen falta un título suficiente (con capacidad, legitimación y escritura pública) y un modo suficiente (la tradición posesoria), que valen lo mismo y se necesitan juntos. Antes de llegar a la escritura, la operación recorre varios instrumentos no obligatorios, que no transmiten el dominio y generan solo derechos personales: el corretaje, las promesas, la reserva, la seña y el refuerzo. Entre ellos, la reserva (un precontrato con suma baja) y la seña (una cláusula accidental, hoy confirmatoria salvo pacto penitencial) se distinguen por su contenido y no por el título que lleven. El boleto de compraventa, que el Código no define y solo regula en sus efectos frente a terceros (arts.~1170 y 1171), divide a la doctrina: para unos es un contrato preliminar (con plazo de caducidad de un año y reclamo limitado a daños), para otros es el contrato definitivo (con posibilidad de exigir la escrituración). Elegir una postura u otra cambia lo que se puede reclamar; por eso el boleto es el punto más importante, y más discutido, de la clase.
\end{important}

\end{document}
